The compact contribution consists of four new descendants of the complete
indexed943/radix67 equations in package~142 \cite{coprime}. The inherited
Pell identities, coprime square extraction, and rational adjacent-Pell
descent are prior work. The new steps are the particular residue and
rational carriers which remove an index quotient, and the smaller modulus
which reduces the degree of one of the residue constructions.

\subsection{The complete polynomial definitions}

The four fixed positive integers $\beta,K_0,M_C,M_F$ are specified in
Appendix~\ref{app:compact} and supplied as complete finite byte strings.
They satisfy
\[
 \beta\ge31,\qquad K_0>0,\qquad 0<M_C<\beta<M_F<2\beta.
\]
These inequalities suffice for some early estimates. The full calibrated
tuple, rather than these inequalities alone, is required for the theorem.
For the first three equations the independent coordinates, all strictly
positive integers, are
\[
 (J,F,\alpha,z,f,n,i,T,s,w,\eta,\zeta,\gamma).
\]
Define the following polynomial abbreviations:
\begin{align*}
 q&=\beta J+1,& C&=q-F-\alpha,& X&=wq,&Y&=sq^3,\\
 E&=XY,&k&=\eta+\zeta,&\ell&=EkY,&c&=kY+\eta,\\
 a&=E+Y,&a_0&=a+2,&H&=4a+3,&\Delta&=a^2+H=a_0^2-1,\\
 D&=X+ac+\gamma H,&A&=1+\ell(\ell+k),
 \end{align*}
and
\begin{align*}
 R&=\bigl[(q-1)(q-F)+\alpha\bigr](q^2-1)+(M_C+qM_F)J,\\
 N_1&=D^2-\Delta c^2,\\
 N_t&=(K_0+w)C+(q-F)-z(q-1).
\end{align*}
The parent has an additional positive coordinate $h$ and factor
$N_k=k-hE-R$. Our new equations omit both. Every letter introduced by a
definition is an abbreviation, not a further quantified unknown.

\paragraph{The 70-operation equation.}
Put
\begin{align*}
 t&=En-k+R,& I&=Ei c^2,\\
 W&=\Delta A I^2,&Q&=\Delta W,&N_s&=f^2-W,\\
 V&=f^2(cT-R)-c,&B&=Q(t^2-AV^2)-t^2+A.
\end{align*}
Then
\begin{equation}\label{eq:p13raw}
 \boxed{P_{13,70}=N_1N_t(3B+1)N_s-1=0.}
\end{equation}
This is a polynomial in exactly thirteen positive unknowns, of exact
total degree 206. Its literal straight-line program uses 70 binary
additions, subtractions and multiplications.

\paragraph{The 71-operation equation.}
Retain $t=En-k+R$, and instead put
\begin{align*}
 m&=AEi,& I&=mc^2,& W&=\Delta I^2,\\
 N_s&=f^2-W,& Q&=\Delta(f^2-1),\\
 V&=f^2(cT-R)-c,&B&=Q(t^2-AV^2)-t^2+A.
\end{align*}
With these definitions use the same expression
\begin{equation}\label{eq:p13reduced}
 \boxed{P_{13,71}=N_1N_t(3B+1)N_s-1=0.}
\end{equation}
Its exact degree is 162, and its complete program has 71 operations.

\paragraph{The smaller-modulus equation.}
In \eqref{eq:p13reduced}, make just the two substitutions
\begin{equation}\label{eq:smallmod}
 t=q^4n-k+R,\qquad m=Aq^2i.
\end{equation}
Call the resulting polynomial $P_{13,72}$. Its exact degree is 154,
and its complete program has 72 operations. The extra operation computes
$q^4$ from the already present $q^2$.

\paragraph{The twelve-variable equation.}
Drop $f$ from the independent coordinate list. Keep all common outer
definitions, and define
\begin{align*}
 m&=a_0Di,& I&=mc^2,&W&=\Delta A I^2,\\
 F_2&=1+W,&Q&=\Delta W,&G&=AF_2,\\
 V&=F_2(cT-R)-c,&e&=k-R-1,&Z&=En+e,\\
 B_{12}&=Q(Z^2-E^2GV^2)-Z^2+E^2G.
\end{align*}
The equation is
\begin{equation}\label{eq:p12}
 \boxed{P_{12,74}=N_1N_t(3B_{12}+1)-1=0.}
\end{equation}
It has twelve positive unknowns, 74 operations and exact degree 438.
There is no rational operation in this definition. Rational quantities
will appear only in its proof.

\subsection{Recovering the first three Pell pairs}

At a zero of any of the three thirteen-variable equations, each factor
of the product in \eqref{eq:p13raw} or \eqref{eq:p13reduced} is an integer
unit. The factor $3B+1$ is 1 modulo 3, so it is $+1$ and $B=0$.
Furthermore $\Delta=a_0^2-1$ is 0 or 3 modulo 4. A norm
$u^2-\Delta v^2$ cannot be $-1$ modulo 4. Thus $N_1=1$.

The exact folded identity is
\begin{equation}\label{eq:fold}
 (Q-1)t^2=A(QV^2-1).
\end{equation}
In the raw construction, $A\mid Q$ and $Q>1$ hold unconditionally.
Equation~\eqref{eq:fold} excludes $t=0$. Both $Q-1$ and $QV^2-1$
are $-1$ modulo $A$. At each prime dividing $A$, valuation in
\eqref{eq:fold} therefore gives
$2v_p(t)=v_p(A)$. Consequently $A=\tau^2$ with $\tau>0$ integral,
and $\tau\mid t$. Only now may we write
$W=\Delta(\tau I)^2$ and exclude $N_s=-1$ modulo 4. Hence $N_s=1$.

In the reduced constructions, $W=\Delta I^2$ from the beginning, so
$N_s=1$ follows first. It gives $Q=(\Delta I)^2>1$. As $A\mid I$,
the same valuation argument again recovers $A=\tau^2$ and $\tau\mid t$.
In all cases $N_t=1$, since it is the only remaining factor.

Put $y=|t|/\tau$. With the positive strong ordinate $I_*$ and $S$ defined by
\[
 I_*=
 \begin{cases}
   \tau Ei c^2,&P_{13,70},\\
   AEi c^2,&P_{13,71},\\
   Aq^2i c^2,&P_{13,72},
 \end{cases}
 \qquad S=\Delta I_*,
\]
we have recovered
\begin{align}
 \tau^2&=1+\ell(\ell+k),& D^2-\Delta c^2&=1,\label{eq:firstmain}\\
 f^2-\Delta I_*^2&=1,&Q&=S^2,\label{eq:strong}\\
 QV^2-(Q-1)y^2&=1,&y&>0.\label{eq:aux}
\end{align}
Both signs of $t$ and $V$ have been allowed. Also $c^2\mid I_*$.

Before interpreting any index, $N_t=1$ gives
\[
 (K_0+w)C=1+z(q-1)-(q-F)\ge F>0.
\]
Thus $C>0$. Set $G_0=(q-1)(q-F)+\alpha=q\alpha+(q-1)C$.
The fixed constant bounds give
\[
 q\le G_0\le q(q-1),\qquad
 0<(M_C+qM_F)J<(2q+1)(q-1),
\]
and therefore
\begin{equation}\label{eq:Rbounds}
 3q<R<q^4-q^3+q^2-1<q^4-2\le E-2.
\end{equation}
In particular none of these bounds has used the omitted $h$.

\subsection{Authenticating the main index without that quotient}

For $b\ge2$, define the Pell sequences by
\[
 \chi_b(0)=1,\quad\chi_b(1)=b,\qquad
 \psi_b(0)=0,\quad\psi_b(1)=1,
 \qquad u_{r+2}=2bu_{r+1}-u_r.
\]
The standard solution classification, addition formulas, monotonicity,
and strong divisibility of $\psi_b$ are retained ordinary lemmas.
For clarity, the particular signed-index argument needed here is included.

Write $v=EY=XY^2$ and $P_0=2v+1$. The first equation in
\eqref{eq:firstmain} is $\tau^2-v(v+1)k^2=1$. Its fundamental
positive solution is $(2v+1,2)$: an ordinate of one would require a
square strictly between $v^2$ and $(v+1)^2$. Hence, for some $r\ge1$,
\[
 \tau=\chi_{P_0}(r),\qquad k=2\psi_{P_0}(r).
\]
The main norm gives $D=\chi_{a_0}(p)$ and $c=\psi_{a_0}(p)$ for
some $p\ge1$. If $r=1$, then $k=2$, $\eta=\zeta=1$, and
$c=2Y+1$ lies strictly between $\psi_{a_0}(1)=1$ and
$\psi_{a_0}(2)=2a_0$. This is impossible. Therefore
\begin{equation}\label{eq:largec}
 r\ge2,\qquad k\ge8XY^2+4,\qquad c>kY>8XY^3>2R.
\end{equation}
Since $P_0>a_0$ and $c>k$, monotonicity also gives $p>r$ and $c>2p$.

The strong norm yields $f=\chi_{a_0}(m)$ and $I_*=\psi_{a_0}(m)$.
From $c^2\mid I_*$, strong divisibility gives $p\mid m$. Writing $m=pu$
and using composition gives
\[
 I_*=c\psi_D(u),\qquad D^2\equiv1\pmod c.
\]
The recurrence implies $\psi_D(u)\equiv uD^{u-1}\pmod c$.
Because $D$ is a unit modulo $c$, one gets $c\mid u$, so
\begin{equation}\label{eq:rank}
 pc\mid m,\qquad c\mid m,\qquad m>2p.
\end{equation}

The auxiliary equation gives an index $\nu\ge1$ with
$S|V|=\chi_S(\nu)$ and $y=\psi_S(\nu)$. It is odd: the even-index
abscissas are $\pm1$ modulo $S$, whereas $S|V|$ is divisible by $S$.
Write $\nu=2u+1$. The integer polynomials
\[
 L_0(Z)=1,\quad L_1(Z)=4Z-3,\quad
 L_{u+1}(Z)=(4Z-2)L_u(Z)-L_{u-1}(Z)
\]
satisfy, by induction,
\begin{align*}
 \chi_b(2u+1)&=bL_u(b^2),&L_u(0)&=(-1)^u(2u+1),\\
 L_u(1-a_0^2)&=(-1)^u\psi_{a_0}(2u+1).
\end{align*}
Here $Q\equiv1-a_0^2\pmod f$, $V\equiv-c\pmod f$,
$Q\equiv0\pmod c$ and $V\equiv-R\pmod c$.
It follows that $\psi_{a_0}(\nu)\equiv\pm c\pmod f$.
Squaring and using $\chi_{a_0}(2j)=1+2\Delta\psi_{a_0}(j)^2$ gives
\[
 \chi_{a_0}(2\nu)\equiv\chi_{a_0}(2p)\pmod{\chi_{a_0}(m)}.
\]

The signed step-down lemma says that, for $0<j<m$, a congruence
$\chi_{a_0}(t)\equiv\chi_{a_0}(j)\pmod{\chi_{a_0}(m)}$
forces $t\equiv\pm j\pmod{2m}$. To see this, use the addition
identities to reduce and reflect the index to a signed
$\chi_{a_0}(s)$, $0\le s\le m$. The endpoint $s=m$ has residue zero.
A negative representative cannot equal the specified positive one,
since $2\chi_{a_0}(m-1)<\chi_{a_0}(m)$; among positive
representatives strict monotonicity forces $s=j$.
Apply it with $j=2p<m$. Then $\nu\equiv\pm p\pmod m$, hence
modulo $c$. Reducing $L_u(Q)$ modulo $c$ separately gives
$\nu\equiv\pm R\pmod c$. Thus $p\equiv\pm R\pmod c$.
The strict inequalities $0<p,R<c/2$ force
\begin{equation}\label{eq:pR}
 p=R.
\end{equation}
This is the point at which the old main index is recovered; the omitted
quotient has not appeared in any premise.

\subsection{How the residue carrier restores the quotient}

For $P_{13,70}$ and $P_{13,71}$, let $M=E$; for $P_{13,72}$, let
$M=q^4$. In each case $t=Mn-k+R$, $P_0\equiv1\pmod M$, and
$M\mid S^2$. The odd-ordinate recurrence gives
\begin{equation}\label{eq:oddres}
 \psi_S(2u+1)\equiv(-1)^u\pmod{4S^2}.
\end{equation}
Also $\tau\equiv1\pmod M$ and $k\equiv2r\pmod M$.
Consequently $t=\pm\tau y\equiv\pm1\pmod M$, and for a sign
$\varepsilon\in\{1,-1\}$,
\[
 2r-R\equiv\varepsilon\pmod M.
\]
Equations~\eqref{eq:largec}--\eqref{eq:pR} give $1\le r<R<M$.
Thus $|2r-R-\varepsilon|<R<M$, forcing the exact equality
$2r=R+\varepsilon$.

The negative sign is impossible. If $R=2r+1$, put $B_0=2a_0^2-1>P_0$.
Composition and monotonicity give
\[
 \psi_{a_0}(2r)=2a_0\psi_{B_0}(r)
     \ge2a_0\psi_{P_0}(r)=a_0k>k(Y+1).
\]
Then $c=\psi_{a_0}(R)>k(Y+1)$, contradicting
$c=kY+\eta<k(Y+1)$. Therefore $2r=R+1$.

The full first-Pell recurrence supplies $k\equiv2r\pmod{4EY}$,
even in the smaller-modulus variant. In every case we can restore
\begin{equation}\label{eq:hrestore}
 h=\frac{k-R-1}{E}\in\mathbb Z_{>0}.
\end{equation}
Positivity follows from $k\ge8XY^2+4>R+1$.
Together with $|t|$ and the appropriate positive rescaled multiplier,
these are a zero of the exact retained parent. Its native/refutation
calibration gives a finite syntactic ZFC refutation.

Conversely, a refutation supplies canonical outer parent witnesses with
$R\equiv3\pmod4$, $A=\tau^2$, $c=\psi_{a_0}(R)$ and $N_k=N_t=N_1=1$.
The finite matrix-period argument permits the strong multiplier to be
chosen divisible by $\tau E$, $AE$, or $Aq^2$, respectively. Choose
fresh auxiliary witnesses for this new strong pair:
\[
 V_*=\frac{\chi_S(R)}{S},\qquad y_*=\psi_S(R),\qquad
 T_*=\frac{V_*+c+Rf^2}{cf^2}.
\]
The odd quotient identities give $V_*\equiv-c\pmod{f^2}$ and
$V_*\equiv-R\pmod c$. Since $\gcd(f,c)=1$ and $f^2\equiv1\pmod c$,
$T_*$ is a positive integer.
Equation~\eqref{eq:oddres} gives $y_*\equiv-1\pmod M$.
Thus $t_*:=\tau y_*$ is $-1$ modulo $M$, while $k-R$ is 1 modulo $M$.
The number $n=(t_*+k-R)/M$ is a positive integer, and every factor
in the new equation is $+1$. This proves the converse existence map.

\subsection{The rational carrier and twelve positive unknowns}

The following inherited descent is included because it is the exact
integrality mechanism required by \eqref{eq:p12}. It is substantially
stronger than merely clearing a denominator in an arbitrary Pell equation.

\begin{lemma}[Adjacent-coefficient integrality]\label{lem:integral}
Let $Q\ge2$ and $V$ be integers. If $y\in\mathbb Q$ satisfies
$QV^2-(Q-1)y^2=1$, then $y\in\mathbb Z$.
\end{lemma}
\begin{proof}
Replace $V,y$ by their absolute values and write $y=b/d>0$ in lowest
terms. Neither coordinate is zero, and $d^2\mid Q-1$. If $V=1$, then
$y=1$. Otherwise $y>V>1$. Put $D_0=Q-1$ and
\[
 V'=(2Q-1)V-2D_0y,\qquad
 y'=\bigl|(2Q-1)y-2QV\bigr|.
\]
The first coordinate is integral. The identity
\[
 ((2Q-1)V)^2-(2D_0y)^2=V^2+4D_0>0
\]
shows $V'>0$, while $y>V$ gives $V'<V$.
Direct expansion preserves the same norm equation. The reduced
denominator remains exactly $d$, since $2Q-1\equiv1\pmod d$.
Descent in the positive integer $V$ reaches $V=1$, where the denominator
is one. Thus the original $d$ was one.
\end{proof}

At a zero of \eqref{eq:p12}, the unit arguments give $N_1=N_t=1$
and $B_{12}=0$. Therefore \eqref{eq:Rbounds} holds. Since
$e=k-R-1\ge1-R>-E$ and $n\ge1$, we have $Z>0$.
The positive rational $t=Z/E=n+e/E$ satisfies
\begin{equation}\label{eq:ratfold}
 (Q-1)t^2=AF_2(QV^2-1).
\end{equation}
Here $A\mid Q$ and $F_2\equiv1\pmod A$. Also
\[
 Q-1=\Delta F_2-a_0^2,
 \qquad QV^2-1\equiv-D^2\pmod{F_2}.
\]
The anchors in $m=a_0Di$ ensure $F_2\equiv1\pmod{a_0}$ and
$F_2\equiv1\pmod D$. Consequently
\[
 \gcd(A,F_2)=1,\qquad
 \gcd\bigl(AF_2,(Q-1)(QV^2-1)\bigr)=1.
\]
Prime valuations in \eqref{eq:ratfold} show that both $A$ and $F_2$
have even valuations. This use of valuations is valid for rational
$t$, whose prime valuations are still integers. Hence
$A=\tau^2$ and $F_2=f^2$ for positive integers $\tau,f$.

Now $Q=(\Delta\tau mc^2)^2>1$, and
$y=t/(\tau f)>0$ satisfies the hypothesis of
Lemma~\ref{lem:integral}. It follows that $y$, and hence $t$, is an
integer. Thus $E\mid e$. Since $e>-E$, one has $e\ge0$.
The possibility $e=0$ would give $k=R+1$. The first Pell equation
forces either $k=2$ or $k\ge8XY^2+4$, both incompatible with
\eqref{eq:Rbounds}. Therefore $h=e/E>0$. These restored $h,t$ give
an exact zero of the anchored thirteen-variable parent, proving soundness.

For completeness, retain the canonical outer data of a refutation,
so $e=hE$ with $h>0$. The parent's matrix-period construction supplies
arbitrarily large strong multipliers divisible by $\tau a_0D$.
Recompute the auxiliary pair after this choice. Its combined root
$t=\tau f y$ can be made larger than $h$. Then $n=t-h>0$, $Z=Et$,
and $B_{12}$ is exactly $E^2$ times the vanishing parent residual.
Thus \eqref{eq:p12} has a positive zero.

\subsection{Exact degrees, operation counts, and logical conclusion}

Deleting the old $N_k$ removes its three operations and one multiplication
in the final product. The first two residue carriers add one multiplier
scaling and three operations for $t$, preserving the parent operation
counts. The third uses a new $q^4$ gate and a lower-degree multiplier.
The twelve-variable formula clears the single rational carrier by $E^2$;
its literal graph has 74 gates. Counts concern the supplied complete
programs, with fixed integer leaves and reused values free.

Exact degree verification accounts for the cancellation
\[
 N_1=2ac(X+\gamma H)+(X+\gamma H)^2-Hc^2,
\]
which makes $\deg N_1=18$. Merely propagating syntactic degrees
through $D^2-\Delta c^2$ would overestimate it. Each output has a
nonzero coefficient at the claimed top degree:
\begin{center}\small
\begin{tabular}{@{}lrl@{}}
\toprule
Equation & Degree & One attained coefficient\\
\midrule
$P_{13,70}$ &206&$-24\beta^{124}$\\
$P_{13,71}$ &162&$-24\beta^{96}$\\
$P_{13,72}$ &154&$-24\beta^{92}$\\
$P_{12,74}$ &438&$+24\beta^{272}$\\
\bottomrule
\end{tabular}
\end{center}
The full monomials, homogeneous leading forms, exact graph checks and
coefficient bindings are included in the verification artifacts.

For each of these four fixed polynomials $P$, the existence maps and
the retained native/refutation calibration prove, in ordinary arithmetic
formalizable in ZFC,
\[
 \mathrm{ZFC}\vdash
 \bigl(\neg\exists\mathbf x\in\mathbb Z_{>0}^{\,d}\;P(\mathbf x)=0\bigr)
 \ \longleftrightarrow\ \Con(\mathrm{ZFC}).
\]
The positive domain is part of the statement. It is not a claim about
unrestricted integer zeros of the same unmodified expression.
